\label{chap:introduction}

\section{Contesto e motivazione}

La scelta di una mossa scacchistica dipende da relazioni che non sono puramente
geometriche: mobilità, attacchi, difese, inchiodature e vincoli di legalità
collegano case anche lontane sulla scacchiera. GraphMate studia se una
rappresentazione esplicita di queste relazioni, elaborata mediante reti neurali a
grafo (\textit{Graph Neural Networks}, GNN), possa sostenere la predizione della
prima mossa risolutiva di un problema tattico di matto.

I puzzle scacchistici costituiscono un caso di studio utile perché associano a una
posizione una soluzione verificabile e impongono regole formali sullo spazio delle
azioni. Il progetto confronta progressivamente una classificazione su un
vocabolario globale di mosse, l'applicazione di vincoli di legalità, lo scoring
diretto delle mosse legali e il reranking basato sulla posizione successiva alla
mossa. In parallelo, valuta una variante che integra un segnale temporale
sintetico, con l'obiettivo di verificare se la specifica informazione di timing
implementata migliori la scelta della mossa.

Come termine di confronto esterno alla famiglia GNN, GraphMate include due modelli
linguistici locali sottoposti a un protocollo controllato di generazione di una
singola mossa UCI. Il confronto non intende misurare ogni possibile capacità
scacchistica dei modelli linguistici: valuta sistemi completi nel compito e nel
formato fissati dal protocollo. La generalizzazione dei modelli viene inoltre
esaminata su composizioni classiche YACPDB, distinte dalla distribuzione dei puzzle
Lichess usata per lo sviluppo principale.

\section{Obiettivi e domande di ricerca}

L'obiettivo generale è costruire e valutare una pipeline riproducibile che trasformi
posizioni scacchistiche in grafi, apprenda a ordinare mosse candidate e renda
tracciabile il passaggio dai dati ai risultati congelati. Gli esperimenti rispondono
a due domande di ricerca:

\begin{description}
    \item[RQ1.] Per quali profondità di matto il modello GNN con timing fornisce
    indicazioni di mossa migliori delle baseline linguistiche selezionate?
    \item[RQ2.] Nell'esperimento controllato, l'integrazione del timing sintetico
    migliora la prestazione rispetto al modello GNN equivalente senza timing?
\end{description}

La prima domanda confronta il Modello B con Qwen 3.5 4B e 9B sul benchmark Lichess,
l'unico per il quale sono disponibili gli ingressi temporali previsti. La seconda
isola invece il contributo del timing attraverso il confronto appaiato tra il
Modello B e A3. Il ruolo dei modelli A, A1, A2 e A4 è ampliare l'analisi delle scelte
architetturali, in particolare dei vincoli di legalità e della valutazione
post-mossa.

\section{Organizzazione della relazione}

Il Capitolo~2 introduce le nozioni scacchistiche e neurali necessarie. I
Capitoli~3 e~4 descrivono rispettivamente i dati e la rappresentazione a grafo; il
Capitolo~5 presenta la famiglia di modelli. Il Capitolo~6 definisce il disegno
sperimentale, mentre il Capitolo~7 riporta le misure. I Capitoli~8 e~9 ne discutono
il significato e formulano le conclusioni. Le appendici raccolgono specifiche
implementative, protocolli, iperparametri e informazioni di riproducibilità.
